\subsection{Setup and evidence reuse}
For $\delta=2$ the required short-gate devices have $L_G=52$~nm, nominal
channel doping and process settings, and $T_{\mathrm{SI}}=30$ or 10~nm.
At $V_{GS}=0$~V and $V_{DS}=1.2$~V, we reuse the archived Part~2
SDevice operating-point TDRs:
\texttt{part2\_potential\_vg00\_vd12\_n53\_des.tdr} (30~nm) and
\texttt{part2\_potential\_vg00\_vd12\_n57\_des.tdr} (10~nm).
No new SProcess or SDevice simulation was necessary.
The raw SVisual exports contain both
\texttt{ConductionBandEnergy} and \texttt{ValenceBandEnergy}.
The longitudinal cutlines follow the lateral $X$ direction
at 1.5~nm inside silicon below the gate-oxide interface,
using $Y=-0.1785\,\mu$m (30~nm) and
$Y=-0.1585\,\mu$m (10~nm).
The vertical cutlines follow the thickness direction at
channel-centre $X=0$. Silicon points were selected by film
coordinates and positive channel acceptor concentration
to exclude oxide-band discontinuities. The vertical graph
depth is 0 at the gate-side surface and 10 or 30~nm
at the respective buried-oxide interface.
The longitudinal originals contain 1628/1633 samples, while
the silicon-only vertical exports contain 37/16.
The repository plot CSVs retain all valid vertical points
and represent the original longitudinal curves using 48/56
originally sampled points, with at most 0.002~eV
piecewise-linear reconstruction error. Physical differences
reported below were audited using the complete original cutlines.

\subsection{Longitudinal energy bands}
\begin{figure}[H]\centering
\begin{tikzpicture}
\begin{axis}[width=.92\linewidth,height=6.8cm,
xlabel={Source-to-drain position, $X$ ($\mu$m)},
ylabel={$E_C$ (eV)},xmin=-.85,xmax=.85,
ymin=-1.5,ymax=.55,grid=both,legend pos=south west]
\addplot+[no marks,thick]
table[col sep=comma,x=X_um,y=Ec_eV]
{data/part5_longitudinal_tsi30.csv};
\addlegendentry{$T_{\mathrm{SI}}=30$ nm}
\addplot+[no marks,thick,dashed]
table[col sep=comma,x=X_um,y=Ec_eV]
{data/part5_longitudinal_tsi10.csv};
\addlegendentry{$T_{\mathrm{SI}}=10$ nm}
\end{axis}
\end{tikzpicture}
\caption{Conduction-band energy, sampled 1.5 nm beneath the
gate-oxide interface at $V_{GS}=0$, $V_{DS}=1.2$ V.
The source lies on the left and the drain on the right.
Shown are genuine simulator points selected for plotting.}
\label{fig:part5_long_ec}
\end{figure}

\begin{figure}[H]\centering
\begin{tikzpicture}
\begin{axis}[width=.92\linewidth,height=6.8cm,
xlabel={Source-to-drain position, $X$ ($\mu$m)},
ylabel={$E_V$ (eV)},xmin=-.85,xmax=.85,
ymin=-2.35,ymax=-.55,grid=both,legend pos=south west]
\addplot+[no marks,thick]
table[col sep=comma,x=X_um,y=Ev_eV]
{data/part5_longitudinal_tsi30.csv};
\addlegendentry{$T_{\mathrm{SI}}=30$ nm}
\addplot+[no marks,thick,dashed]
table[col sep=comma,x=X_um,y=Ev_eV]
{data/part5_longitudinal_tsi10.csv};
\addlegendentry{$T_{\mathrm{SI}}=10$ nm}
\end{axis}
\end{tikzpicture}
\caption{Valence-band energy on the same two longitudinal
cutlines and at identical terminal biases.}
\label{fig:part5_long_ev}
\end{figure}

\paragraph{Source-channel barrier.}
With the source conduction band referenced at
$X\simeq-0.8\,\mu$m, $E_C$ equals $-0.16156$ eV
(30 nm) or $-0.15609$ eV (10 nm).
The maximum $E_C$ along each line occurs at
$X=-0.1275\,\mu$m and $-0.1343\,\mu$m,
respectively, with values $0.39918$ and
$0.38485$ eV. Hence the extracted local
source-to-crest barrier heights are
$0.56074$ and $0.54094$ eV. In particular,
the thin-film cutline barrier is about
20 meV \emph{lower} than the thick-film
one under these specific bias conditions.
It would be incorrect to claim a taller
thin-film band-edge crest from these results,
even though the measured thin-film
$I_{\mathrm{OFF}}$ is lower (Part~2).
Leakage depends on the complete two-dimensional
transport path rather than one horizontal
surface-adjacent line.

\paragraph{Drain-side slope and electrostatic lowering.}
From $X\simeq0.20$ to $0.40\,\mu$m,
$E_C$ falls by 13.9 meV for 30 nm
and by 7.2 meV for 10 nm, before
a steeper descent through the drain junction.
This particular pre-junction band-edge slope
is smaller for the thinner film.
The source/drain $E_C$ plateaus differ by
approximately 1.20 eV, consistent with
$V_{DS}=1.2$ V. Drain electrostatic coupling
can lower the source-channel barrier in short
devices, but \emph{quantitative drain-induced
barrier lowering} requires a corresponding
low-$V_{DS}$ band snapshot at the same gate
bias. The present two TDRs both use
$V_{DS}=1.2$ V, so this study does not
invent a numerical change in barrier height
between drain biases.

\subsection{Channel-centre vertical energy bands}
\begin{figure}[H]\centering
\begin{tikzpicture}
\begin{axis}[width=.92\linewidth,height=6.8cm,
xlabel={Depth from gate-side silicon interface (nm)},
ylabel={$E_C$ (eV)},xmin=0,xmax=30,
ymin=.375,ymax=.47,grid=both,legend pos=north west]
\addplot+[no marks,thick]
table[col sep=comma,x=Depth_nm,y=Ec_eV]
{data/part5_vertical_tsi30.csv};
\addlegendentry{$T_{\mathrm{SI}}=30$ nm}
\addplot+[no marks,thick,dashed]
table[col sep=comma,x=Depth_nm,y=Ec_eV]
{data/part5_vertical_tsi10.csv};
\addlegendentry{$T_{\mathrm{SI}}=10$ nm}
\end{axis}
\end{tikzpicture}
\caption{Channel-centre vertical conduction-band energy
at $V_{GS}=0$ V and $V_{DS}=1.2$ V.
The 10-nm trace ends at the 10-nm Si/BOX interface;
the 30-nm trace ends at 30 nm. Only silicon is shown.}
\label{fig:part5_vert_ec}
\end{figure}

\begin{figure}[H]\centering
\begin{tikzpicture}
\begin{axis}[width=.92\linewidth,height=6.8cm,
xlabel={Depth from gate-side silicon interface (nm)},
ylabel={$E_V$ (eV)},xmin=0,xmax=30,
ymin=-.735,ymax=-.635,grid=both,legend pos=north west]
\addplot+[no marks,thick]
table[col sep=comma,x=Depth_nm,y=Ev_eV]
{data/part5_vertical_tsi30.csv};
\addlegendentry{$T_{\mathrm{SI}}=30$ nm}
\addplot+[no marks,thick,dashed]
table[col sep=comma,x=Depth_nm,y=Ev_eV]
{data/part5_vertical_tsi10.csv};
\addlegendentry{$T_{\mathrm{SI}}=10$ nm}
\end{axis}
\end{tikzpicture}
\caption{Channel-centre vertical valence-band energy,
using the same off-gate/high-drain operating points
as Fig.~\ref{fig:part5_vert_ec}.}
\label{fig:part5_vert_ev}
\end{figure}

\paragraph{Vertical bending and gate control.}
The 30-nm device exhibits an increase in
$E_C$ from 0.39347 eV near the gate to
0.45855 eV near the buried oxide:
a total of 65.1 meV through the top
silicon. The corresponding increase for
10-nm silicon is 0.38269 to 0.39275 eV,
or just 10.1 meV.
The valence bands show similar changes
of 64.9 meV and 10.4 meV, respectively.
Thus the thinner silicon film displays a more
nearly uniform transverse electrostatic band
profile, indicating reduced variation of
potential across the body. This is consistent
with improved geometric gate control and
the 10-nm device's reduced OFF current
observed in Part~2. Both energy bands
shift together within each film, with a local
silicon bandgap of about 1.105 eV.

\paragraph{Fully depleted versus partially depleted.}
The original vertical SDevice records, before
removing auxiliary columns for the report plots,
contain mobile-hole densities far below the
local acceptor densities throughout both films.
The maximum ratios of hole density to acceptor
concentration were $7.5\times10^{-8}$
(30 nm) and $5.8\times10^{-8}$ (10 nm).
The recorded net space charge was also
approximately the negative acceptor density
across the entire channel-centre silicon,
within about 0.01\% in these samples.
There is therefore no observed neutral,
hole-rich core along either cutline:
\emph{both devices are consistent with full
body depletion at $V_{GS}=0$ V,
$V_{DS}=1.2$ V}. The larger 30-nm band
bending is not by itself evidence of
partial depletion. The given data do not
establish full depletion at all terminal
biases, but do directly support depleted
bodies at the specified Part~5 condition.

\paragraph{Evidence limitations.}
The archived TDRs were used without rerunning
the physics. The four figures are based on
Sentaurus numerical samples, not idealized
band sketches. The longitudinal plot
subsamples reproduce the original raw
traces to within 2 meV under linear
interpolation; vertical silicon samples
were not subsampled. These band-edge
snapshots cannot independently produce a
numerical DIBL value without low-drain-bias
band profiles; any stronger claim would
exceed the evidence obtained here.

